% !TeX root = ../papers/nc-p3.tex
\section{四元群次数を持つ混合交換子のトーラスと全閉点の正則性}
\label{sec:marked-klein-mixed-torus}
\begin{theorem}\label{thm:marked-klein-mixed-torus}
二つの三次元空間の座標軸に三つの非自明な四元群文字を対応させ、
四組 $Q_i,P_i,G_i,F_i$ の各第 $i$ 行列が $E_{jk},E_{kj}$
（$i,j,k$ は巡回）の二成分だけを持つとする。
これら二十四成分が全て非零で $[Q_i,P_i]\ne0$ である印付き窓を全て考える。
その二切断方程式の正規形は下記の三変数主枝と一次元の付加枝で尽きる。
主枝の全次数Koszul左右AS正則性は一つの数 $u^6\ne1$ と同値である。
途中でWalsh正規化の対角分母が消える全ての点もこの判定に含まれる。
付加枝は前節の中心一次拡大で常に正則であり、二枝が一致する値も明示できる。
正則点は全て周期が2を割る。印を忘れた全同型と最小周期は、
孤立した単純印の有限比較、軸置換とトーラスの単項式方程式で必要十分に判定できる。
\end{theorem}
\begin{proof}[四組の二成分を全て解く]
$E_{jk}$ の係数を正、$E_{kj}$ の係数を負と呼ぶ。
$\delta_i=Q_i^+P_i^--Q_i^-P_i^+$ とすれば跡条件 $\sum[Q_i,P_i]=0$ より
三つとも同じ $\delta\ne0$ で、対応する $G,F$ の差もこの $\delta$ である。
対角の残余基底と全体倍率を用いて
\[
 Q_i^+=u,\quad Q_i^-=1,\quad P_i^+=v_i r_i,\quad P_i^-=r_i=(u-v_i)^{-1}
\]
とする。実際 $Q_i^-$ は添字の三倍率で揃えられ、$Q_i^+/Q_i^-$ は
三座標の隣接比の平方で変えるため、積を固定して共通値 $u$ にできる。倍率で $\delta=1$ に戻す。
いずれの $u,v_i$ も非零で、$v_i\ne u$ である。

非対角式は二つの階数一行列 $G_i^+ F_j^+$ と $G_i^- F_j^-$ の対角外六成分を指定する。
$G_i^+=c_i,G_i^-=e_i$ とすれば、巡回の順で
\[
 F_i^+=-r_{i-1}/c_{i-1},\quad F_i^-=-r_i/e_{i-1},\qquad
 c_{i+1}/c_{i-1}=u v_i r_i/r_{i-1},\quad
 e_{i+1}/e_{i-1}=u v_{i+1}r_{i+1}/r_i .
\]
この系の必要十分条件は $u^3 v_0 v_1 v_2=1$。$c_i$ は共通倍率を除き一意、
$e_i=h c_{i+1}$ で、$c_0$ は残余のスカラーで1にできる。
$G_i^+ F_i^- -G_i^-F_i^+=1$ はちょうど
\[
                        (hu+1)(h v_i-1)=0\qquad (i=0,1,2)
\]
となる。従って主枝では $h=-1/u$、付加選択は全ての $v_i=v$ が等しく $h=1/v$ のときだけである。
このとき $uv$ は三乗根で、$u,v_i$ を同じ三乗根倍にする残余を用い $uv=1$ にできる。
この規格化では $u^2=-1$ でのみ二根が合流する。
$Q_i,P_i$ の交叉交換子はそこでの規格化 $uv=1$ の全体で、非対角が交代行列となり、
前節の $t=(u+1)/(u-1)$ に当たる（$u=1,-1$ は $v\ne u$ に反する）。
主選択は $G,F$ が $W(t)$ 側、付加選択は一方だけを替えた中心一次側である。
これは $G_i^-/G_{i+1}^+$ の値を代入しても分かる。
\end{proof}

\begin{proof}[Walsh正規化と極限のAS判定]
主枝で四ラベルを $0,1,2,3\in(\mathbf Z/2)^2$ とする（この段では非零ラベルを $1,2,3$ と読む）。
隣接ラベルが等しい成分は零である。残る三十六成分
$w_{i,i+b,i+b+e,i+e}$, $b,e\ne0$ で、四つの相異なるラベルなら元の非零な一星成分である。
二ラベルが反復する成分は $w_{0i0i}=w_{i0i0}=1$、
$w_{jkjk}=w_{kjkj}=s_i$ であり、
\[
       s_i=u(r_j+r_k)-1=(u^2-v_j v_k)/((u-v_j)(u-v_k)) .
\]
まず全 $s_i\ne0$ とする。四位置の四つずつの対角倍率と九個の非零数 $m_{be}$ に関する
単項式系を作り、各 $w$ を $(-1)^{i_1(b+e)_0}m_{be}$ に変える。
指数行列 $M$ の行は、四位置の該当ラベルに1、該当する $m$ に $-1$ を持つ。
三十六行二十五列のSmith標準形の非零対角は $1$ が十四、$2$ が五、$4$ が二である。
$\operatorname{sign}/w$ を右辺とした十五の左核積は全て1なので、代数閉体上で全点に解がある。
ここで検算は独立不定元 $u,a=v_0,b'=v_1$、$v_2=(u^3ab')^{-1}$ における有理恒等式である。
$N$ をWalshの符号付き平方行列とすると、$N$ の符号は行・列別の積なので次の比には残らない。
\[
 \frac{N_{11}N_{22}}{N_{12}N_{21}}-1
       =\frac{(u^6-1)(ab'u^4-1)}{u^4(u-a)(u-b')}.
\]
具体的に $i,b,e$ の辞書式順（$i$ は0、$b,e$ は1始まり）の指数行のうち、
零始まりで $9,10,12,13,18,19,21,22$ を順に $-1,1,1,-1,-1,1,1,-1$ 倍して
足すと、求める平方比の指数行になる。
\path{computations/verify_klein_mixed_torus.py} が全八十一式、整数左核の全積とこの式を検算する。
$ab'u^4-1=u/v_2-1\ne0$。三色ずつあり、行和・列和が零なので全二次余因子は符号まで等しい。
従って定理\ref{thm:walsh-normal-quadrics}により判定は $u^6\ne1$ となる。

$s_i=0$ が一つでもあれば $v_i=u^{-5}\ne u$、従って既に $u^6\ne1$ である。
全位置でラベル $0,i$ を同時に $z$ 倍し、テンソルを $z^2$ で割る。
この操作を欠けた全 $i$ について合成すると四つの異なるラベルの成分は不変であり、
対角色 $i$ の $0,i$ の対はさらに $z^2$ 倍、残る対は既に零である。他の対角色は保たれる。
極限では $s_i=0$ の各色の対角 $m_{ii}$ を丸ごと零とした配置になる。
残った単項式系の整数左核は、消した行に零を補えば元の整数左核に入る。
残る右辺はそこで非零であるから左核の恒等式も特殊化でき、極限をWalsh表示にできる。
各切断は階数6以下で、非零な三つの色ブロックは各々少なくとも二つの非零係数を持つ。
Walshブロックの左右二次行列表現により非零ブロックの階数は少なくとも2
（可逆な外部符号変更の後で二次行列による正則表現である）。
従って三ブロックとも階数2である。$N$ の非対角は全て非零で一つ以上の対角が零なので
階数1にはなれない。極限で全次数正則となり、
命題\ref{prop:unrestricted-good-reduction}を $k[[z]]$ 上で適用すると元の窓も全次数正則である。
\end{proof}

\begin{proof}[全同型と周期の有限判定]
$\delta\ne0$ なので、$\langle Q,P\rangle$ と $\langle G,F\rangle$ は各々
非対角行列の六空間そのものである。その跡直交補は各側の跡零対角Cartan部分代数で、
三つの共通固有直線が軸枠を復元する。印を保つ同型は従って両側の単項行列に限る。
添字と残る二軸の対応から軸置換は同じ置換であり、残る倍率は四組の係数の単項式比較で決まる。
例えば $v_i$ が全て等しい場合以外では主枝しかなく、巡回置換と同じ三乗根による
$(u,v_i)$ の一斉乗法、または奇置換と四数全ての逆数を取って同じ三乗根を掛ける操作が全てである。
実際 $Q_i^+$ と $Q_i^-$ の比に作用する三つの倍率は積1で、
共通の $u$ を共通値に送れば同じ三乗根となる。逆に平方根の対角比を選び
$Q_i^-$ と $\delta$ の規格化を戻せばこの変換は全て実現する。
一致した $v_i$ の二枝の比較は前節の明示解、または四組の単項式比較で行える。

一位置移動の $Q',P'$ は $F,G$ である。これを再規格化すると主枝では
$u'=-1/u,v'_i=-1/v_i$ を取れる。実際 $F_i^+/F_i^-=h c_i r_{i-1}/(c_{i-1}r_i)$
の積は $h^3$、$u'=h$ に揃えた後の $v_i'$ は $1/(hu v_i)$ である。
非定数な三つ組では二移動で元に戻り、定数の場合も前節の周期を用いればよい。
一般に印を忘れるには、基準の単純な星点を固定し相手の同位置の全ての孤立単純点対を調べる。
これらは有限個であり、点対応全体に正次元成分があってもよい。
星点対は局所一次式が六独立なので常に孤立単純である。対応する他の二点は
反対評価の階数3の核から求め、適合しないものを除く。
跡零規格化の後、上記の六空間・Cartanとその固有軸を検査し、それぞれの枠で
全四組の係数を比較する。これは添字保持同型の必要十分条件である。
一移動との同じ有限比較で最小周期1か2かを判定できる。
\end{proof}
